\section{Related Work}

Our work builds on and extends three lines of research: benchmark generalization studies, underspecification in machine learning, and texture bias analysis.

\subsection{Benchmark Generalization Studies}

\citet{recht2019imagenet} provided the first systematic evidence that ImageNet classifiers do not generalize to ImageNet-like test sets. By creating ImageNetV2 with a carefully matched distribution, they demonstrated consistent 10--15\% accuracy drops across 70+ models---a finding that challenged assumptions about benchmark reliability. Subsequent work extended this analysis to CIFAR-10 with the CINIC-10 dataset \citep{darlow2018cinic}, finding similar generalization gaps. However, these studies examine individual datasets in isolation without connecting generalization failure to dataset popularity or the broader ML ecosystem's optimization patterns.

\citet{engstrom2020identifying} and \citet{taori2020measuring} further documented ``effective robustness'' differences across model families, showing that ImageNet accuracy improvements do not uniformly translate to out-of-distribution performance. Our work extends this line by proposing that \emph{popularity itself} is a predictor of generalization failure---not merely an artifact of individual dataset characteristics.

\subsection{Underspecification and Benchmark Overfitting}

\citet{damour2020underspecification} introduced the underspecification framework, demonstrating that models achieving identical benchmark performance can exhibit dramatically different behavior under distribution shift. Their theoretical analysis explains \emph{why} benchmark performance fails to predict deployment success: the optimization surface admits many equivalent solutions, and benchmarks do not constrain which one the model learns. This provides theoretical grounding for our hypothesis: if popular benchmarks receive more optimization pressure, models may converge to increasingly benchmark-specific solutions.

The benchmark co-evolution concept also relates to Goodhart's Law in ML \citep{thomas2022reliance}---when a measure becomes a target, it ceases to be a good measure. We operationalize this concern by measuring the actual research investment differential between high-use and low-use datasets.

\subsection{Texture Bias and Feature Learning}

\citet{geirhos2019imagenet} demonstrated that ImageNet-trained CNNs rely heavily on texture rather than shape, contrary to human perception. \citet{hermann2020origins} traced this bias to training data statistics, suggesting that benchmark-specific artifacts shape learned representations. We hypothesized that intensive optimization on popular benchmarks might amplify texture bias as a mechanism for the co-evolution effect.

However, our experiments contradict this expectation at CIFAR scale: ResNet-18 shows \emph{lower} texture bias (0.126) than VGG-11 (0.161). This suggests that architectural innovations (skip connections, batch normalization) may counteract texture exploitation, and that ImageNet-scale findings do not directly transfer to 32$\times$32 image classification. Our negative result narrows the mechanism search: the co-evolution effect, while real, operates through pathways other than texture bias.

\subsection{Dataset Documentation and Ecosystem Analysis}

\citet{gebru2021datasheets} proposed Datasheets for Datasets to standardize documentation, addressing concerns about unexamined data practices. While their work focuses on documentation quality, our study examines dataset \emph{usage} patterns---specifically, how uneven research attention across the benchmark ecosystem correlates with generalization outcomes. These perspectives are complementary: both highlight that benchmark selection and maintenance practices deserve more scrutiny than they currently receive.
